240
5 Differential Calculus
Summing these inequalities over $i$ from $1$ to $n$, we obtain
$$ \frac{\sum_{i=1}^n x_i y_i}{X^{1/p}Y^{1/q}} \le 1, $$
which is equivalent to relation (5.88).
We obtain (5.89) similarly from (5.87). Since equality occurs in (5.86) and (5.87)
only when $a = b$, we conclude that it is possible in (5.88) and (5.89) only when a
proportionality $x_i^p = \lambda y_i^q$ or $y_i^q = \lambda x_i^p$ holds.
c. Minkowski's Inequalities$^{18}$
Let $x_i \ge 0, y_i \ge 0, i = 1, \dots, n$. Then
$$ \left(\sum_{i=1}^n (x_i + y_i)^p\right)^{1/p} \le \left(\sum_{i=1}^n x_i^p\right)^{1/p} + \left(\sum_{i=1}^n y_i^p\right)^{1/p} \quad \text{when } p > 1, \quad (5.90) $$
and
$$ \left(\sum_{i=1}^n (x_i + y_i)^p\right)^{1/p} \ge \left(\sum_{i=1}^n x_i^p\right)^{1/p} + \left(\sum_{i=1}^n y_i^p\right)^{1/p} \quad \text{when } p < 1, p \ne 0. \quad (5.91) $$
Proof We apply Hölder's inequality to the terms on the right-hand side of the iden-
tity
$$ \sum_{i=1}^n (x_i + y_i)^p = \sum_{i=1}^n x_i(x_i + y_i)^{p-1} + \sum_{i=1}^n y_i(x_i + y_i)^{p-1}. $$
The left-hand side is then bounded from above (for $p > 1$) or below (for $p < 1$)
in accordance with inequalities (5.88) and (5.89) by the quantity
$$ \left(\sum_{i=1}^n x_i^p\right)^{1/p} \left(\sum_{i=1}^n (x_i + y_i)^q\right)^{1/q} + \left(\sum_{i=1}^n y_i^p\right)^{1/p} \left(\sum_{i=1}^n (x_i + y_i)^q\right)^{1/q} . $$
After dividing these inequalities by $\left(\sum_{i=1}^n (x_i + y_i)^p\right)^{1/q}$, we arrive at (5.90) and
(5.91).
Knowing the conditions for equality in Hölder's inequalities, we verify that
equality is possible in Minkowski's inequalities only when the vectors $(x_1,\dots,x_n)$
and $(y_1,\dots, y_n)$ are collinear.
For $n = 3$ and $p = 2$, Minkowski's inequality (5.90) is obviously the triangle
inequality in three-dimensional Euclidean space.
$^{18}$H. Minkowski (1864–1909) – German mathematician who proposed a mathematical model
adapted to the special theory of relativity (a space with a sign-indefinite metric).